\documentclass[11pt,letterpaper]{article}

\usepackage[spanish,es-noshorthands]{babel}
\usepackage{fontspec}
\setmainfont{Source Serif Pro}
\usepackage{microtype}
\usepackage{geometry}
\usepackage{booktabs}
\usepackage{tabularx}
\usepackage{enumitem}
\usepackage{xcolor}
\usepackage{graphicx}
\usepackage{csquotes}
\usepackage{amsmath}
\usepackage{siunitx}
\usepackage[version=4]{mhchem}
\usepackage{tikz}
\usetikzlibrary{arrows.meta,positioning,shapes.geometric,fit,backgrounds,patterns,decorations.pathreplacing}
\usepackage{xurl}
\usepackage[hidelinks]{hyperref}
\usepackage[backend=biber,style=apa,sorting=nyt]{biblatex}
\DeclareLanguageMapping{spanish}{spanish-apa}

\geometry{margin=2.3cm}
% Los rótulos de las figuras usan nodos de ancho fijo con texto alineado, y la
% bibliografía compone en modo sloppy: ambos generan avisos de línea holgada
% que son cosméticos. Se silencian para que `make check` siga siendo sensible
% a los desbordes reales.
\hbadness=10000
\setlength{\emergencystretch}{1em}
\setlength{\parindent}{0pt}
\setlength{\parskip}{0.42em}
\setlist{nosep}
\renewcommand{\topfraction}{0.92}
\renewcommand{\bottomfraction}{0.85}
\renewcommand{\textfraction}{0.06}
\renewcommand{\floatpagefraction}{0.72}
\setcounter{topnumber}{2}
\setcounter{totalnumber}{3}
\addbibresource{referencias.bib}
\AtBeginBibliography{\sloppy\small}
\setlength{\bibitemsep}{0.28em}
\setlength{\bibhang}{1.4em}

\sisetup{
  output-decimal-marker = {,},
  group-separator = {.},
  group-minimum-digits = 4,
  detect-all
}

% Colores: un color por paradigma de comunicación (banda S, banda X, banda Ka
% y óptico), el rojo de los límites físicos, el gris de los ejes, el violeta
% del caso de contraste y el verde de lo que se sostiene.
\definecolor{cS}{RGB}{62,132,168}
\definecolor{cX}{RGB}{26,58,102}
\definecolor{cKa}{RGB}{186,86,30}
\definecolor{cOpt}{RGB}{42,124,78}
\definecolor{cLim}{RGB}{176,42,42}
\definecolor{cGris}{RGB}{128,124,118}
\definecolor{cNet}{RGB}{104,62,148}
\definecolor{cVerde}{RGB}{56,120,86}
\definecolor{cDato}{RGB}{36,102,148}

% --- Figura 1: la envolvente de la Red del Espacio Profundo -----------------
% x: año 1955 a 2035; y: logaritmo decimal de bits por segundo a 5 UA
\newcommand{\ex}[1]{((#1)-1955)*0.1800}
\newcommand{\ey}[1]{((#1)+13.6)*0.2450}
% barras horizontales de tasa (paneles B de las figuras 1 y 3)
\newcommand{\tx}[1]{(#1)*4.30}
% #1 y, #2 tasa, #3 color, #4 rótulo izquierdo, #5 rótulo derecho
\newcommand{\barraT}[5]{%
  \node[left, text width=4.55cm, align=right, font=\scriptsize] at (-0.12,#1)
       {#4};%
  \fill[#3] (0,#1-0.185) rectangle ({\tx{#2}},#1+0.185);%
  \node[right, font=\scriptsize\bfseries, text=#3] at ({\tx{#2}+0.12},#1)
       {#5};%
}
% la misma barra con la escala ampliada que necesita la figura 3
\newcommand{\ux}[1]{(#1)*21.00}
\newcommand{\barraU}[5]{%
  \node[left, text width=4.55cm, align=right, font=\scriptsize] at (-0.12,#1)
       {#4};%
  \fill[#3] (0,#1-0.20) rectangle ({\ux{#2}},#1+0.20);%
  \node[left, font=\scriptsize\bfseries, text=white] at ({\ux{#2}-0.12},#1)
       {#5};%
}

% --- Figura 2: las asíntotas y la escalera de frecuencia --------------------
\newcommand{\dx}[1]{(#1)*3.55}
% #1 y, #2 décadas extraídas, #3 décadas restantes, #4 rótulo, #5 nota
\newcommand{\barraA}[5]{%
  \node[left, text width=3.9cm, align=right, font=\scriptsize] at (-0.12,#1)
       {#4};%
  \fill[cDato] (0,#1-0.185) rectangle ({\dx{#2}},#1+0.185);%
  \draw[cLim, pattern=north east lines, pattern color=cLim!55]
       ({\dx{#2}},#1-0.185) rectangle ({\dx{#2}+\dx{#3}},#1+0.185);%
  \node[right, font=\scriptsize, text=cGris!40!black, align=left]
       at ({\dx{#2}+\dx{#3}+0.14},#1) {#5};%
}
\newcommand{\fy}[1]{(#1)*0.455}

% --- Figura 3: la velocidad de acceso a internet ----------------------------
\newcommand{\nx}[1]{((#1)-1983)*0.2980}
\newcommand{\ny}[1]{((#1)-2.0)*0.5150}

% --- Figura 4: las tres series y los tres niveles ---------------------------
\newcommand{\rx}[1]{((#1)-1955)*0.1560}
\newcommand{\ry}[1]{(#1)*4.30}

\title{\vspace{-1.2cm}\textbf{Veinte órdenes de magnitud y cuatro asíntotas}\\
\large Las cuatro afirmaciones sobre el progreso tecnológico, medidas sobre
los datos de la Red del Espacio Profundo y contrastadas con la velocidad de
acceso a internet}
\author{Carlos Luis Rojas Aragonés\\
\small Gestión del Desarrollo Tecnológico: Estrategias y Análisis de Portfolio}
\date{\small 22 de septiembre de 2026}

\begin{document}

\maketitle
\vspace{-0.8em}

\section{Las cuatro afirmaciones, respondidas en corto}

El planteamiento ofrece cuatro afirmaciones sobre cómo progresan las
tecnologías y pide decidir cuáles son verdaderas. Mi respuesta corta es que la
pregunta, tal como está redactada, no tiene una respuesta única, y que eso no
es una evasiva sino el resultado de medir: \textbf{las cuatro afirmaciones son
la misma curva mirada a tres alturas distintas}, y cada una es verdadera a una
altura y falsa a las otras. El Cuadro~\ref{cuadro:ficha} lo resume; el resto
del documento lo sostiene con los números del propio caso.

\begin{table}[b]
\small
\caption{Las cuatro afirmaciones del planteamiento, con el veredicto y el
apartado donde se sostiene con el cálculo. \enquote{Décadas al año} es el
número de órdenes de magnitud que la figura de mérito gana cada año.}
\label{cuadro:ficha}
\begin{tabularx}{\textwidth}{@{}p{0.4cm}>{\raggedright\arraybackslash}p{3.35cm}>{\raggedright\arraybackslash}p{2.5cm}X@{}}
\toprule
& \textbf{Afirmación} & \textbf{Veredicto} & \textbf{Por qué} \\
\midrule
1 & Todas saturan en una asíntota; el progreso exponencial no existe
  & Cierta la premisa,\newline falsa la conclusión
  & Las asíntotas existen y aquí se \emph{demuestran}, no se observan: en 1995
    la codificación de la DSN ya había consumido el 80 por ciento de toda la
    ganancia que el teorema de Shannon permitirá jamás. Pero la función ganó
    20,8 órdenes de magnitud en 64 años (apartado~\ref{sec:af1}) \\
\addlinespace[2pt]
2 & Exponencial con tasa decreciente por el ingreso marginal decreciente
  & Cierta aquí,\newline falsa como ley
  & La envolvente de la DSN cae de 0,889 a 0,121 décadas al año, así que la
    conclusión se cumple. Pero el acceso a internet sostuvo 0,179 durante 42
    años, y el mecanismo que actúa no es un rendimiento marginal genérico sino
    asíntotas con nombre (apartado~\ref{sec:af2}) \\
\addlinespace[2pt]
3 & Exponencial a tasa constante por alternancia de curvas S
  & \textbf{La más útil}\newline de las cuatro
  & El mecanismo es literalmente lo que dibuja el gráfico de la JPL: banda S,
    banda X, banda Ka y óptico. Pero la alternancia \emph{sostiene} el
    progreso, no garantiza la tasa: en la DSN alterna y aun así decelera
    (apartado~\ref{sec:af3}) \\
\addlinespace[2pt]
4 & Exponencial con tasa creciente por las capacidades humanas
  & No la sostiene\newline ningún dato realizado
  & El único tramo acelerado del gráfico de la JPL es el que todavía no ha
    ocurrido. Y \textcite{bloom2020} miden lo contrario: sostener la tasa de
    Moore exige hoy 18 veces más investigadores que en 1971
    (apartado~\ref{sec:af4}) \\
\bottomrule
\end{tabularx}
\end{table}

\subsection*{Lo que encontré, dicho por adelantado}

\begin{quote}
\textbf{La afirmación 3 describe el mecanismo, la 1 describe el componente, y
la 2 y la 4 son propiedades empíricas de cada tecnología concreta, no leyes.}
La Red del Espacio Profundo ganó \textbf{20,8 órdenes de magnitud en 64 años},
lo que refuta la afirmación 1 en su conclusión; pero su tasa se derrumbó de
0,889 décadas al año en la era de banda S a 0,121 en la de banda Ka, con una
\textbf{meseta plana de doce años entre 1993 y 2005}, lo que confirma la
afirmación 2 y refuta la 4. La velocidad de acceso doméstico a internet, en
cambio, sostuvo \textbf{0,179 décadas al año durante 42 años}, un
\num{51} por ciento anual, atravesando tres medios físicos distintos: eso
confirma la afirmación 3 en su enunciado completo. Misma pregunta, dos
respuestas opuestas, medidas con la misma unidad.
\end{quote}

\subsection*{Nota sobre las cifras}

Distingo tres clases de número, porque mezclarlas es la forma más rápida de
que un argumento sobre curvas S parezca más sólido de lo que es.

\textbf{Lecturas del gráfico} de la capacidad normalizada de la JPL
\parencite{descanso2015}, que es el dato primario de todo el apartado~3. Son
lecturas visuales de un eje logarítmico de 24 décadas: les asigno una
tolerancia de \textbf{$\pm$0,2 décadas} por punto. Eso basta para las
pendientes, que se calculan sobre diferencias de 2 a 16 décadas, y no basta
para ningún escalón individual menor de medio orden de magnitud.

\textbf{Valores citados} de fuentes técnicas: la temperatura de ruido del
sistema de 20,9 K y la relación $E_b/N_0$ de 0,65 dB conseguida por Galileo
vienen de \textcite{layland1997}; la temperatura del fondo cósmico de
microondas, 2,725 K, de \textcite{fixsen2009}; las velocidades de los módems y
de las redes de acceso, de las recomendaciones de la Unión Internacional de
Telecomunicaciones citadas en el apartado~\ref{sec:internet}.

\textbf{Cálculos propios}, declarados donde se usan: las fronteras entre eras,
las pendientes en décadas al año, las ganancias teóricas por cambio de banda,
la capacidad de Shannon del canal telefónico y el reparto entre ganancia ya
extraída y margen restante del Cuadro~\ref{cuadro:asintotas}.

\section{El criterio que decide: a qué altura se mide}

Antes de los datos hace falta el criterio, porque sin él las cuatro
afirmaciones son indecidibles. La disciplina que propone \textcite{deweck2022}
es declarar la figura de mérito sobre \emph{la función prestada} y no sobre la
implementación, y es exactamente ahí donde las cuatro afirmaciones se separan.
Hay tres alturas posibles:

\begin{itemize}
\item \textbf{Componente}: una pieza concreta con una física concreta. El
  máser refrigerado, el reflector de 70 metros, el código convolucional de
  longitud 7. Aquí la figura de mérito tiene un techo que se puede escribir
  como una ecuación.
\item \textbf{Paradigma}: un conjunto coherente de componentes que comparten
  un principio de funcionamiento. La banda S, la banda X, la banda Ka, el
  enlace óptico. Un paradigma agota sus componentes y muere.
\item \textbf{Función}: lo que el usuario quiere, con independencia de cómo se
  consiga. Aquí: \emph{mover bits desde una sonda lejana hasta la Tierra}.
\end{itemize}

La figura de mérito funcional de este caso es la que usa la propia JPL en su
gráfico de capacidad: \textbf{la tasa de datos del canal, en bits por segundo,
normalizada a una distancia de 5 unidades astronómicas}, es decir, a la
distancia de Júpiter. La normalización importa porque la potencia recibida cae
con el cuadrado de la distancia, de modo que comparar la tasa bruta de una
sonda lunar con la de una sonda a Neptuno no mide tecnología, mide geometría.

Con esa figura de mérito declarada, las cuatro afirmaciones dejan de ser
opiniones. La 1 es una afirmación sobre componentes. La 3 es una afirmación
sobre cómo se encadenan los paradigmas. Y la 2 y la 4 son afirmaciones sobre
la \emph{pendiente de la envolvente funcional}, que es un número que se mide.
El Cuadro~\ref{cuadro:eras} y la Figura~\ref{fig:dsn} lo miden.

\section{Los datos: veinte órdenes de magnitud en sesenta y cuatro años}
\label{sec:datos}

\begin{figure}[!htb]
\centering
\begin{tikzpicture}

% ================= Panel A: la envolvente ==============================
\node[anchor=south west, font=\small\bfseries] at (-0.95,6.22)
     {Panel A. La capacidad de la Red del Espacio Profundo, 1958 a 2033};
\node[anchor=south west, font=\scriptsize\itshape, text=cGris] at (-0.95,5.92)
     {bits por segundo a 5 unidades astronómicas, escala logarítmica};

% rejilla y ejes
\foreach \v in {-12,-9,-6,-3,0,3,6,9} {
  \draw[cGris!25] (0,{\ey{\v}}) -- (14.45,{\ey{\v}});
  \node[left, font=\scriptsize, text=cGris] at (-0.07,{\ey{\v}})
       {$10^{\v}$};
}
\draw[cGris!70] (0,0) -- (0,5.72);
\draw[cGris!70] (0,0) -- (14.45,0);
\foreach \a in {1960,1970,1980,1990,2000,2010,2020,2030} {
  \draw[cGris!70] ({\ex{\a}},0) -- ({\ex{\a}},-0.09);
  \node[below, font=\scriptsize, text=cGris] at ({\ex{\a}},-0.09) {\a};
}
\node[below, font=\scriptsize, text=cGris] at (7.2,-0.50) {año};

% --- era de banda S ---
\draw[cS, very thick]
  ({\ex{1958}},{\ey{-13}}) -- ({\ex{1960}},{\ey{-13}}) --
  ({\ex{1960}},{\ey{-6.0}}) -- ({\ex{1961.5}},{\ey{-6.0}}) --
  ({\ex{1961.5}},{\ey{-5.5}}) -- ({\ex{1963}},{\ey{-5.5}}) --
  ({\ex{1963}},{\ey{-4.0}}) -- ({\ex{1965}},{\ey{-4.0}}) --
  ({\ex{1965}},{\ey{-2.7}}) -- ({\ex{1966}},{\ey{-2.7}}) --
  ({\ex{1966}},{\ey{-1.3}}) -- ({\ex{1969}},{\ey{-1.3}}) --
  ({\ex{1969}},{\ey{-0.5}}) -- ({\ex{1971}},{\ey{-0.5}}) --
  ({\ex{1971}},{\ey{0.5}})  -- ({\ex{1972}},{\ey{0.5}}) --
  ({\ex{1972}},{\ey{1.0}})  -- ({\ex{1973}},{\ey{1.0}}) --
  ({\ex{1973}},{\ey{2.0}})  -- ({\ex{1974}},{\ey{2.0}}) --
  ({\ex{1974}},{\ey{2.6}})  -- ({\ex{1975}},{\ey{2.6}}) --
  ({\ex{1975}},{\ey{3.0}})  -- ({\ex{1976}},{\ey{3.0}});

% --- era de banda X ---
\draw[cX, very thick]
  ({\ex{1976}},{\ey{3.0}})  -- ({\ex{1976}},{\ey{4.3}}) --
  ({\ex{1977}},{\ey{4.3}})  -- ({\ex{1977}},{\ey{4.7}}) --
  ({\ex{1979}},{\ey{4.7}})  -- ({\ex{1979}},{\ey{5.0}}) --
  ({\ex{1981}},{\ey{5.0}})  -- ({\ex{1981}},{\ey{5.15}}) --
  ({\ex{1986}},{\ey{5.15}}) -- ({\ex{1986}},{\ey{5.3}}) --
  ({\ex{1987}},{\ey{5.3}})  -- ({\ex{1987}},{\ey{5.4}}) --
  ({\ex{1989}},{\ey{5.4}})  -- ({\ex{1989}},{\ey{5.5}}) --
  ({\ex{1993}},{\ey{5.5}})  -- ({\ex{1993}},{\ey{5.7}}) --
  ({\ex{2005}},{\ey{5.7}});

% --- era de banda Ka ---
\draw[cKa, very thick]
  ({\ex{2005}},{\ey{5.7}})  -- ({\ex{2006}},{\ey{5.7}}) --
  ({\ex{2006}},{\ey{6.0}})  -- ({\ex{2017}},{\ey{6.0}}) --
  ({\ex{2017}},{\ey{7.05}}) -- ({\ex{2022}},{\ey{7.05}}) --
  ({\ex{2022}},{\ey{7.78}}) -- ({\ex{2026}},{\ey{7.78}});

% --- era óptica, proyectada ---
\draw[cOpt, very thick, dash pattern=on 4pt off 2pt]
  ({\ex{2013}},{\ey{2.0}})  -- ({\ex{2022}},{\ey{2.0}}) --
  ({\ex{2022}},{\ey{5.7}})  -- ({\ex{2023}},{\ey{5.7}}) --
  ({\ex{2023}},{\ey{6.8}})  -- ({\ex{2025}},{\ey{6.8}}) --
  ({\ex{2025}},{\ey{7.6}})  -- ({\ex{2027}},{\ey{7.6}}) --
  ({\ex{2027}},{\ey{8.0}})  -- ({\ex{2030}},{\ey{8.0}}) --
  ({\ex{2030}},{\ey{8.6}})  -- ({\ex{2033}},{\ey{8.6}});

% --- rectas de tendencia por era ---
\draw[cGris!85, dotted, thick]
  ({\ex{1958}},{\ey{-13}}) -- ({\ex{1976}},{\ey{3.0}});
\draw[cGris!85, dotted, thick]
  ({\ex{1977}},{\ey{4.3}}) -- ({\ex{1993}},{\ey{5.7}});
\draw[cGris!85, dotted, thick]
  ({\ex{2006}},{\ey{6.0}}) -- ({\ex{2022}},{\ey{7.78}});

% --- anotaciones ---
\node[right, font=\scriptsize, text=cS, align=left, inner sep=3pt]
     at ({\ex{1958.4}},{\ey{-12.6}}) {Explorer 1 (1958)};
\node[right, font=\scriptsize, text=cS, align=left, inner sep=3pt]
     at ({\ex{1966.5}},{\ey{-9.6}}) {\textbf{banda S:} 0,889 déc./año};
\draw[cS!70, thin] ({\ex{1966.2}},{\ey{-4.1}}) -- ({\ex{1969.4}},{\ey{-2.6}});
\node[right, font=\scriptsize, text=cS, align=left, inner sep=3pt]
     at ({\ex{1969.3}},{\ey{-4.6}}) {antena de 64 m (G), 1966};
\node[right, font=\scriptsize, text=cX, align=left, inner sep=3pt]
     at ({\ex{1976.4}},{\ey{3.20}}) {paso a \textbf{banda X} (1976)};
\node[above, font=\scriptsize, text=cX, align=center, inner sep=3pt]
     at ({\ex{1983}},{\ey{5.80}}) {\textbf{banda X:} 0,076 déc./año};
\draw[cLim, thick] ({\ex{1993}},{\ey{5.7}}) -- ({\ex{2005}},{\ey{5.7}});
\node[below, font=\scriptsize\bfseries, text=cLim, align=center, inner sep=4pt]
     at ({\ex{1997.5}},{\ey{5.50}}) {meseta de 12 años\\0,007 déc./año};
\draw[cKa!75, thin] ({\ex{2006.4}},{\ey{5.90}}) -- ({\ex{2008.3}},{\ey{4.30}});
\node[right, font=\scriptsize, text=cKa, align=left, inner sep=3pt]
     at ({\ex{2008.1}},{\ey{3.85}}) {\textbf{banda Ka} (2006)\\0,121 déc./año};
\draw[cKa!75, thin] ({\ex{2018.0}},{\ey{6.90}}) -- ({\ex{2020.0}},{\ey{5.70}});
\node[left, font=\scriptsize, text=cKa, align=right, inner sep=3pt]
     at ({\ex{2033}},{\ey{5.35}}) {conjuntos de 34 m (G)};
\node[left, font=\scriptsize, text=cOpt, align=right, inner sep=3pt]
     at ({\ex{2012.6}},{\ey{2.0}}) {óptico: LLCD (2013)};
\node[left, font=\scriptsize, text=cOpt, align=right, inner sep=3pt]
     at ({\ex{2016.4}},{\ey{8.70}}) {\textbf{óptico (proyectado):}
     0,228 déc./año};

% ================= Panel B: la tasa por era =============================
\begin{scope}[xshift=5.60cm, yshift=-5.90cm]
\node[anchor=south west, font=\small\bfseries] at (-4.75,3.55)
     {Panel B. La tasa de progreso de la envolvente, era por era};

\barraT{3.00}{0.889}{cS}{\textbf{Banda S}, 1958 a 1976\\16,0 décadas en 18
  años}{0,889}
\barraT{2.25}{0.076}{cX}{\textbf{Banda X}, 1977 a 1993\\1,22 décadas en 16
  años}{0,076}
\barraT{1.50}{0.007}{cLim}{\textbf{Meseta}, 1993 a 2005\\0,08 décadas en 12
  años}{0,007}
\barraT{0.75}{0.121}{cKa}{\textbf{Banda Ka}, 2006 a 2022\\1,93 décadas en 16
  años}{0,121}
\barraT{0.00}{0.228}{cOpt}{\textbf{Óptico}, 2023 a 2031\\\emph{proyectado},
  1,82 décadas en 8 años}{0,228}

\draw[cGris!70] (0,-0.32) -- ({\tx{0.95}},-0.32);
\foreach \t/\r in {0/0, 0.2/{0,2}, 0.4/{0,4}, 0.6/{0,6}, 0.8/{0,8}} {
  \draw[cGris!70] ({\tx{\t}},-0.32) -- ({\tx{\t}},-0.41);
  \node[below, font=\scriptsize, text=cGris] at ({\tx{\t}},-0.41) {\r};
}
\node[below, font=\scriptsize, text=cGris] at ({\tx{0.45}},-0.78)
     {décadas de mejora por año};
\end{scope}

\end{tikzpicture}
\caption{Panel A: la capacidad normalizada de la Red del Espacio Profundo,
redibujada a partir del gráfico de escalones de la JPL
\parencite{descanso2015}. Cada escalón es una tecnología que entra en
servicio; el color marca el paradigma de radiofrecuencia. Las rectas punteadas
son la tendencia de cada era. La lectura clave no es el recorrido total, que
es enorme, sino que \emph{las rectas punteadas no son paralelas}: cada era
progresa más despacio que la anterior. Panel B: esa misma información como
tasa. La primera era avanzó siete veces más deprisa que la de banda Ka, y
entre 1993 y 2005 la envolvente no se movió. Elaboración propia sobre
\textcite{descanso2015}.}
\label{fig:dsn}
\end{figure}

La Figura~\ref{fig:dsn} redibuja el gráfico de escalones que la JPL publica
como medida oficial de la capacidad de la red \parencite{descanso2015}, y el
Cuadro~\ref{cuadro:eras} recoge la aritmética.

El recorrido total es el dato que todo el mundo cita: de $10^{-13}$ bits por
segundo a la distancia de Júpiter con Explorer 1 en 1958 a unos $6\times10^{7}$
con los conjuntos de siete antenas de 34 metros en banda Ka en 2022. Son
\textbf{20,8 órdenes de magnitud en 64 años}, una media de 0,325 décadas al
año, es decir, multiplicar por 2,11 cada año durante seis décadas. Ninguna
tecnología terrestre de uso civil se acerca a eso.

Y esa media es \textbf{engañosa}, que es el hallazgo del apartado. Está
dominada por los primeros dieciocho años: si se retira la era de banda S, los
46 años restantes suman 4,8 décadas, es decir 0,104 décadas al año, veinte
veces menos.

\begin{table}[t]
\small
\caption{La envolvente de la Red del Espacio Profundo, era por era. Los
extremos son lecturas del gráfico de la JPL con una tolerancia de $\pm$0,2
décadas; las pendientes y los tiempos de duplicación son cálculo propio. La
última fila no ha ocurrido: es la proyección que la propia JPL publica.}
\label{cuadro:eras}
\begin{tabularx}{\textwidth}{@{}>{\raggedright\arraybackslash}p{2.35cm}>{\raggedright\arraybackslash}p{1.8cm}>{\raggedleft\arraybackslash}p{1.35cm}>{\raggedleft\arraybackslash}p{1.5cm}>{\raggedleft\arraybackslash}p{1.65cm}X@{}}
\toprule
\textbf{Era} & \textbf{Años} & \textbf{Décadas} & \textbf{Déc./año} &
\textbf{Duplica en} & \textbf{Qué la cerró} \\
\midrule
Banda S & 1958--1976 & 16,00 & \textbf{0,889} & 4,1 meses
  & La banda de 2,3 GHz agotó su ganancia por apertura y por potencia \\
\addlinespace[2pt]
Banda X & 1977--1993 & 1,22 & 0,076 & 47,3 meses
  & Se agotaron el reflector de 70 m, el máser y la codificación concatenada \\
\addlinespace[2pt]
Meseta & 1993--2005 & 0,08 & \textbf{0,007} & 45,6 años
  & Nada la cerró: no hubo paradigma nuevo ni presupuesto para uno \\
\addlinespace[2pt]
Banda Ka & 2006--2022 & 1,93 & 0,121 & 29,9 meses
  & La absorción atmosférica impide subir más en radiofrecuencia \\
\addlinespace[2pt]
\emph{Óptico (proy.)} & \emph{2023--2031} & \emph{1,82} & \emph{0,228} &
  \emph{15,8 meses} & \emph{Sin cerrar; el salto a 193 THz abre 7,6 décadas
  teóricas} \\
\midrule
\textbf{Total realizado} & \textbf{1958--2022} & \textbf{20,78} &
  \textbf{0,325} & \textbf{11,1 meses} & \\
\textbf{Sin la era de S} & \textbf{1977--2022} & \textbf{4,78} &
  \textbf{0,104} & \textbf{34,8 meses} & \\
\bottomrule
\end{tabularx}
\end{table}

\section{Afirmación 1: las asíntotas existen, y aquí se demuestran}
\label{sec:af1}

\textbf{Enunciado:} \emph{todas las tecnologías terminan alcanzando un punto de
saturación, lo que dificulta su progreso más allá de una asíntota determinada.
El progreso exponencial no existe.}

Mi respuesta es que la primera frase es cierta y la segunda es falsa, y que el
caso de la DSN es el mejor sitio posible para verlo porque aquí las asíntotas
no se conjeturan a partir de los datos: \textbf{se deducen de la física y
luego se comprueba que los datos las respetan}. Eso es infrecuente. En la
mayoría de las tecnologías la asíntota se postula porque la curva se aplana, y
eso es un argumento circular. Aquí no.

El Cuadro~\ref{cuadro:asintotas} recoge las cuatro palancas que la DSN
controla desde tierra, cada una con su asíntota, el origen físico de esa
asíntota y cuánta ganancia queda por extraer. El Panel A de la
Figura~\ref{fig:asintotas} lo dibuja.

\begin{table}[t]
\small
\caption{Las cuatro palancas de tierra, su asíntota y el margen que queda. La
columna de la derecha es la conclusión incómoda: en tres de las cuatro
palancas, la mayor parte de la ganancia que la física permitirá jamás ya está
gastada. Cálculo propio sobre los valores citados.}
\label{cuadro:asintotas}
\begin{tabularx}{\textwidth}{@{}p{2.4cm}X>{\raggedright\arraybackslash}p{2.6cm}>{\raggedleft\arraybackslash}p{1.65cm}@{}}
\toprule
\textbf{Palanca} & \textbf{La asíntota y de dónde sale} &
\textbf{Extraído y lo que queda} & \textbf{Gastado} \\
\midrule
\textbf{Codificación}
& El teorema de \textcite{shannon1948} fija el mínimo absoluto en
  $E_b/N_0 = -1{,}59$ dB. Sin codificar hacen falta unos 9,6 dB; Galileo, con
  un código concatenado $(14,1/4)$ y decodificación en cuatro etapas, bajó a
  0,65 dB \parencite{layland1997}
& 8,95 dB / 2,24 dB\newline (0,89 / 0,22 déc.)
& \textbf{80,0 \%} \\
\addlinespace[2pt]
\textbf{Temperatura de ruido}
& El suelo es el fondo cósmico de microondas, 2,725 K \parencite{fixsen2009},
  más la atmósfera. La red opera hoy a 20,9 K en banda X sobre la antena de
  70 m en el cénit \parencite{layland1997}
& hasta 20,9 K /\newline factor 7,7\newline (0,88 déc.)
& \textbf{67,9 \%} \\
\addlinespace[2pt]
\textbf{Apertura}
& La ganancia va con el cuadrado del diámetro, pero un reflector orientable se
  deforma por su propio peso. El paso de 26 a 64 m dio 0,78 décadas; el de 64
  a 70 m, en 1988, dio 0,08
& 0,86 déc. /\newline no hay más\newline en un solo plato
& \textbf{$\approx$100 \%} \\
\addlinespace[2pt]
\textbf{Banda}
& La ganancia va con el cuadrado de la frecuencia. S a X da 1,13 décadas; X a
  Ka da 1,16. Por encima de 32 GHz la absorción del vapor de agua hace
  inutilizable la ventana
& 2,29 déc. /\newline 0 dentro de\newline radiofrecuencia
& \textbf{$\approx$100 \%} \\
\bottomrule
\end{tabularx}
\end{table}

Los tres números que más me importan de ese cuadro son estos:

\begin{enumerate}
\item \textbf{La codificación estaba al 80 por ciento de su límite absoluto en
  1995.} Los ingenieros de la JPL lo sabían y lo escribieron: en 1985
  publicaron un artículo titulado \enquote{En busca de una ganancia de
  codificación de 2 dB} \parencite[véase la bibliografía
  de][]{layland1997}. Ese título es el retrato de una asíntota: cuando una
  disciplina titula sus artículos con la cifra exacta de lo que le queda por
  ganar, la afirmación 1 se cumple sin discusión.
\item \textbf{La mejora de la antena de 64 a 70 metros, en 1988, valió 0,08
  décadas}, un 20 por ciento de capacidad. Fue una obra civil de años
  \parencite{layland1997}. Ese es el aspecto exacto del final de una curva S:
  coste creciente, ganancia decreciente.
\item \textbf{La escalera de bandas de radiofrecuencia está terminada.} S a X
  y X a Ka dieron 1,13 y 1,16 décadas cada una, casi idénticas, porque la
  ganancia va con $f^2$ y las dos razones de frecuencia son parecidas. El
  siguiente peldaño no existe dentro de la radiofrecuencia.
\end{enumerate}

Y a pesar de todo eso, la función ganó 20,8 órdenes de magnitud. \textbf{Por
eso la segunda frase de la afirmación 1 es falsa}: el progreso exponencial no
solo existe, sino que este caso es uno de los ejemplos más extremos que se
conocen. Lo que no existe es progreso exponencial \emph{indefinido en un
componente}, que es una afirmación mucho más modesta y que nadie discute.

Aquí quiero recoger y matizar mi propia intuición inicial, que era que en la
DSN \enquote{las limitaciones eran más regulaciones humanas que otra cosa},
por lo que recuerdo de la entrevista al doctor Deutsch \parencite{deutsch2026}.
Sigo pensando que eso es cierto \emph{hoy}, y el apartado~\ref{sec:af4} lo
desarrolla. Pero al hacer las cuentas he tenido que corregirme: las
limitaciones de la codificación y de la temperatura de ruido no son
regulatorias en absoluto. Son el teorema de Shannon y el fondo cósmico de
microondas, y no admiten negociación con nadie.

\begin{figure}[!htb]
\centering
\begin{tikzpicture}

% ================= Panel A: las cuatro asíntotas =========================
\node[anchor=south west, font=\small\bfseries] at (-4.05,3.05)
     {Panel A. Cuánta ganancia queda en cada palanca de tierra};

\barraA{2.45}{0.89}{0.22}{\textbf{Codificación}\\límite de Shannon}
  {80 \% consumido}
\barraA{1.75}{1.86}{0.88}{\textbf{Temperatura de ruido}\\fondo cósmico, 2,725 K}
  {68 \%}
\barraA{1.05}{0.86}{0.03}{\textbf{Apertura de un plato}\\26 m a 70 m}
  {$\approx$100 \%}
\barraA{0.35}{2.29}{0.02}{\textbf{Banda de radio}\\S a X a Ka}{$\approx$100 \%}

\draw[cGris!70] (0,0.05) -- ({\dx{3.0}},0.05);
\foreach \t/\r in {0/0, 0.5/{0,5}, 1/1, 1.5/{1,5}, 2/2, 2.5/{2,5}, 3/3} {
  \draw[cGris!70] ({\dx{\t}},0.05) -- ({\dx{\t}},-0.04);
  \node[below, font=\scriptsize, text=cGris] at ({\dx{\t}},-0.04) {\r};
}
\node[below, font=\scriptsize, text=cGris] at ({\dx{1.5}},-0.40)
     {décadas de mejora};
\fill[cDato] (0.90,-1.28) rectangle (1.30,-1.01);
\node[right, font=\scriptsize, text=cGris!30!black] at (1.40,-1.145)
     {ya extraído};
\draw[cLim, pattern=north east lines, pattern color=cLim!55]
     (4.20,-1.28) rectangle (4.60,-1.01);
\node[right, font=\scriptsize, text=cGris!30!black] at (4.70,-1.145)
     {queda hasta el límite físico};

% ================= Panel B: la escalera de frecuencia ====================
\begin{scope}[yshift=-6.85cm]
\node[anchor=south west, font=\small\bfseries] at (-0.55,4.28)
     {Panel B. Por qué el siguiente peldaño no está en radiofrecuencia};
\node[anchor=south west, font=\scriptsize\itshape, text=cGris] at (-0.55,3.98)
     {décadas de ganancia teórica del salto, por la dependencia con $f^2$};

\foreach \v in {0,2,4,6} {
  \draw[cGris!25] (0,{\fy{\v}}) -- (12.6,{\fy{\v}});
  \node[left, font=\scriptsize, text=cGris] at (-0.07,{\fy{\v}}) {\v};
}
\draw[cGris!70] (0,0) -- (0,{\fy{7.95}});
\draw[cGris!70] (0,0) -- (12.6,0);

\fill[cX] (1.10,0) rectangle (2.50,{\fy{1.13}});
\node[above, font=\scriptsize\bfseries, text=cX] at (1.80,{\fy{1.13}})
     {1,13};
\node[below, font=\scriptsize, text=cX, align=center] at (1.80,-0.10)
     {S a X\\2,3 a 8,4 GHz\\\emph{1976}};

\fill[cKa] (4.10,0) rectangle (5.50,{\fy{1.16}});
\node[above, font=\scriptsize\bfseries, text=cKa] at (4.80,{\fy{1.16}})
     {1,16};
\node[below, font=\scriptsize, text=cKa, align=center] at (4.80,-0.10)
     {X a Ka\\8,4 a 32 GHz\\\emph{2006}};

\fill[cLim!25] (7.10,0) rectangle (8.50,{\fy{0.001}});
\draw[cLim, thick] (7.10,0.02) -- (8.50,0.02);
\node[above, font=\scriptsize\bfseries, text=cLim, align=center]
     at (7.80,{\fy{0.30}}) {no existe};
\node[below, font=\scriptsize, text=cLim, align=center] at (7.80,-0.10)
     {Ka a algo mayor\\el vapor de agua\\cierra la ventana};

\fill[cOpt] (10.10,0) rectangle (11.50,{\fy{7.56}});
\node[above, font=\scriptsize\bfseries, text=cOpt] at (10.80,{\fy{7.56}})
     {7,56};
\node[below, font=\scriptsize, text=cOpt, align=center] at (10.80,-0.10)
     {Ka a óptico\\32 GHz a 193 THz\\\emph{2023}};

\node[anchor=north west, font=\scriptsize\itshape, text=cGris!30!black,
      text width=12.2cm, align=left] at (-0.50,-1.28)
     {Los dos primeros saltos valen casi lo mismo porque las razones de
      frecuencia son casi iguales. El tercer peldaño no existe: por encima de
      32 GHz la atmósfera absorbe. Por eso el paradigma siguiente no es una
      banda de radio mejor sino un láser, y por eso ofrece de golpe más
      ganancia teórica que todo lo ganado en radiofrecuencia desde 1958.};
\end{scope}

\end{tikzpicture}
\caption{Panel A: las cuatro palancas que la red controla desde tierra, con la
ganancia ya extraída y la que queda hasta el límite físico. Ninguna de las
cuatro barras rayadas es larga, y dos son prácticamente inexistentes: eso es
la afirmación 1, comprobada. Panel B: la escalera de frecuencia, que es la
razón por la que la afirmación 3 puede seguir cumpliéndose pese al Panel A.
Elaboración propia sobre \textcite{layland1997}, \textcite{shannon1948} y
\textcite{fixsen2009}.}
\label{fig:asintotas}
\end{figure}

\section{Afirmación 2: cierta para esta envolvente, y no por la razón que dice}
\label{sec:af2}

\textbf{Enunciado:} \emph{las tecnologías experimentan un progreso exponencial,
pero la tasa de progreso tiende a disminuir con el paso del tiempo, porque
todas se ven afectadas por la ley de la disminución del ingreso marginal.}

Esta era mi respuesta inicial y sigo pensando que su conclusión es correcta
para este caso. Los datos del Cuadro~\ref{cuadro:eras} la confirman sin
ambigüedad: \textbf{0,889, luego 0,076, luego 0,007, luego 0,121 décadas al
año}. La secuencia es decreciente salvo el repunte final, y ese repunte no
recupera ni la séptima parte de la tasa original.

Pero hay dos correcciones que tuve que hacerme a mí mismo al medir, y las dos
importan más que el veredicto.

\textbf{Primera: la palabra \enquote{todas} no se sostiene.} Es una
generalización que el propio caso no puede probar, porque un caso no prueba
una ley universal, y que además tiene contraejemplos medidos. El
apartado~\ref{sec:internet} presenta uno, y no es el único: \textcite{koh2006}
estudian el progreso funcional de las tecnologías de la información durante
más de un siglo y encuentran tasas exponenciales sostenidas, no decrecientes,
en varias funciones; \textcite{farmer2016} y \textcite{nagy2013} llegan a lo
mismo analizando decenas de series de coste y rendimiento. Si la disminución
del ingreso marginal fuera una ley universal, esas series no existirían.

\textbf{Segunda, y esta es la que más me hizo cambiar de opinión: el mecanismo
que enuncia la afirmación no es el que actúa aquí.} \enquote{Ingreso marginal
decreciente} describe una situación en la que cada unidad adicional de esfuerzo
rinde un poco menos que la anterior, de forma suave y continua. Lo que la
Figura~\ref{fig:dsn} muestra no es eso. Muestra tres cosas distintas:

\begin{itemize}
\item \textbf{Un orden de extracción.} Las primeras dieciséis décadas salieron
  de poner una antena direccional y un transmisor de verdad en una nave que
  antes llevaba una antena de látigo. No fue ingenio marginal: fue recoger lo
  que estaba en el suelo. Cualquier tecnología nueva tiene un tramo así, y
  confundirlo con una tasa sostenible es el error de previsión más común.
\item \textbf{Asíntotas duras y nombrables}, las del
  Cuadro~\ref{cuadro:asintotas}, que no producen una pendiente suavemente
  decreciente sino un tope.
\item \textbf{Una meseta de doce años que no tiene causa técnica.} Entre 1993 y
  2005 la envolvente se movió 0,08 décadas. No fue que costara más conseguir
  la mejora siguiente: fue que la mejora siguiente exigía cambiar de
  paradigma, y cambiar de paradigma en una red de antenas de decenas de metros
  es una decisión de presupuesto plurianual, no de ingeniería. Volveré sobre
  esto en el apartado~\ref{sec:af4}, porque es donde mi intuición inicial
  sobre las \enquote{regulaciones humanas} sí acierta de lleno.
\end{itemize}

Dicho de otro modo: la afirmación 2 acierta en el resultado y se equivoca en la
causa, y en gestión tecnológica la causa es lo que sirve para decidir. Si uno
cree que la desaceleración es un ingreso marginal decreciente, la respuesta es
invertir más en lo mismo. Si uno sabe que es una asíntota física más una
decisión de capital aplazada, la respuesta es preparar el paradigma siguiente,
que es lo que la JPL acabó haciendo con el enlace óptico.

\section{Afirmación 3: es el mecanismo correcto, pero no garantiza la tasa}
\label{sec:af3}

\textbf{Enunciado:} \emph{las tecnologías experimentan un progreso exponencial
y, en algunos casos, como el de las tecnologías funcionales, la tasa de
progreso se mantiene más o menos constante, por la alternancia entre curvas
S.}

Esta es, de las cuatro, \textbf{la más útil}, y comparto la intuición que ya
tenía: cuando una curva S tiende a la baja, viene otra tecnología a
reemplazarla. La quiero afinar en dos direcciones, una a favor y otra en
contra.

\textbf{A favor, y es un argumento fuerte:} el gráfico de la JPL no es que sea
compatible con la alternancia de curvas S, es que \emph{la dibuja
explícitamente}. Los cuatro colores de la Figura~\ref{fig:dsn} son cuatro
paradigmas, y el gráfico original de la JPL incluso lo advierte en su pie: los
saltos grandes de capacidad proceden de los cambios de banda de frecuencia
\parencite{descanso2015}. La secuencia es limpia: la banda S agota su apertura
y su potencia hacia 1975, y en 1976 entra la banda X con un salto de 1,3
décadas; la banda X agota el máser, el reflector de 70 metros y la
codificación concatenada hacia 1993, y en 2006 entra la banda Ka; la banda Ka
agota la ventana atmosférica, y en 2023 entra el óptico, que ya ha demostrado
267 megabits por segundo desde 0,19 unidades astronómicas
\parencite{dsoc2023}. Cada paradigma muere por asíntotas propias, que son las
del Cuadro~\ref{cuadro:asintotas}, y su sucesor no lo mejora: lo sustituye. Es exactamente lo que describe
\textcite{foster1986} y lo que \textcite{sood2005} encuentran al examinar
empíricamente la evolución de varias tecnologías: la unidad de análisis útil
no es la curva sino la familia de curvas que se relevan.

\textbf{En contra, y es la parte que no había visto antes de medir:} la
alternancia \emph{sostiene} el progreso, pero no lo mantiene a tasa constante.
La DSN alterna cuatro veces y aun así su tasa cae un factor de siete. La razón
es que un relevo de paradigma no es gratuito ni instantáneo: entre que la
banda X se agota, hacia 1993, y que la banda Ka entra en servicio, en 2006,
pasan doce años en los que no ocurre nada. \textbf{El hueco entre curvas S es
tan parte del fenómeno como el salto.} Y como la afirmación 3 solo menciona el
salto, describe bien el mecanismo y mal la consecuencia.

Donde la afirmación 3 se cumple entera, incluida la parte de la tasa constante,
es en el caso del apartado siguiente. Y creo que la diferencia entre los dos
casos es la lección más transferible de todo el ejercicio.

\section{Afirmación 4: no la sostiene ningún dato realizado de este caso}
\label{sec:af4}

\textbf{Enunciado:} \emph{las tecnologías experimentan un progreso exponencial
y la tasa de progreso aumenta con el paso del tiempo como resultado del
desarrollo de las capacidades de la humanidad.}

Mi respuesta inicial fue que esto es cierto, porque todo eso ayuda al
desarrollo. Al medir he tenido que separar dos cosas que había juntado:
\textbf{que una capacidad ayude no significa que la tasa aumente}. Son
afirmaciones distintas y solo la segunda es la que hay que comprobar.

Sobre este caso la respuesta es limpia y negativa. \textbf{Ningún tramo
realizado de la Figura~\ref{fig:dsn} es más rápido que el tramo anterior,
salvo el paso de la meseta a la banda Ka, que es una recuperación parcial y no
una aceleración.} El único tramo del gráfico cuya pendiente supera a la
anterior de forma clara, el óptico a 0,228 décadas al año, es \emph{el que
todavía no ha ocurrido}. Y conviene notar que es exactamente el tramo donde
más barata resulta el optimismo.

Hay además un argumento cuantitativo general en contra, y es el que más peso
tiene: \textcite{bloom2020} miden la productividad de la investigación en
varios sectores y encuentran que \textbf{mantener la duplicación de densidad
de transistores de la ley de Moore exige hoy más de 18 veces los
investigadores que hacían falta a comienzos de los años setenta}. Esa cifra es
la refutación directa de la afirmación 4 tal como está enunciada: las
capacidades humanas agregadas sí han crecido, y el efecto observable no ha
sido una tasa creciente sino una tasa constante comprada cada vez más cara.

Dicho lo cual, la afirmación 4 tiene un caso real a su favor, y sería
deshonesto no ponerlo. \textcite{sevilla2022} miden el cómputo de entrenamiento
de los sistemas de aprendizaje automático y encuentran un cambio de régimen
nítido: antes de 2010 se duplicaba cada 21,3 meses, en línea con la ley de
Moore, es decir 0,170 décadas al año; desde 2010 se duplica cada 5,7 meses, es
decir \textbf{0,634 décadas al año}, casi cuatro veces más rápido. Eso es una
aceleración medida, no proyectada. Pero es una aceleración de \emph{un insumo}
en \emph{un dominio} durante \emph{una década}, y convertir eso en una ley
general del progreso tecnológico, que es lo que hace \textcite{kurzweil2001},
es el mismo error de extrapolación que cometería quien tomase los 0,889
décadas al año de la era de banda S como la tasa normal de la DSN.

\subsection*{Dónde sí acierta la intuición sobre las limitaciones humanas}

Queda una pieza que creo que es la más interesante del caso, y es donde mi
lectura inicial de la entrevista al doctor Deutsch \parencite{deutsch2026} sí
da en el blanco, aunque no como argumento a favor de la afirmación 4 sino como
matiz a las afirmaciones 2 y 3.

\textbf{La restricción activa de la DSN se ha desplazado de la física a las
instituciones.} En 1965 lo que limitaba la capacidad era el ruido del
receptor. En 2026 lo que la limita es, por este orden: el reparto
internacional del espectro, que decide qué bandas puede usar una sonda; el
ciclo presupuestario plurianual, que decide cuándo se construye una antena; la
edad del parque, con reflectores instalados en los años sesenta que hay que
mantener mientras se construyen los nuevos; y las ventanas de lanzamiento, que
fijan cuándo una tecnología nueva puede volar. Ninguna de esas cuatro cosas es
un problema de ingeniería, y las cuatro juntas explican la meseta de 1993 a
2005 mucho mejor que cualquier asíntota.

Para un director de tecnología eso tiene una consecuencia práctica inmediata:
en un sistema maduro, \textbf{el calendario del relevo entre curvas S lo fija
la organización, no el laboratorio}, y por tanto el hueco entre curvas es una
variable de gestión y no un dato de la naturaleza.

\section{Un caso distinto: la velocidad de acceso a internet desde 1984}
\label{sec:internet}

\begin{quote}
El planteamiento pide exponer brevemente un caso tecnológico distinto que me
interese y cuyo sistema haya evolucionado con el tiempo. El mío es \textbf{la
velocidad de conexión a internet desde un domicilio}, desde los módems de
finales de los ochenta hasta la fibra de hoy. Me interesa porque es el caso
espejo del de la DSN: la misma unidad de medida, el mismo tipo de gráfico y la
misma alternancia de paradigmas, pero \emph{el resultado opuesto}. Entre 1984
y 2026 la velocidad pasó de 300 bits por segundo a 10 gigabits por segundo,
7,52 órdenes de magnitud, atravesando tres medios físicos completamente
distintos: el canal telefónico de voz, el par de cobre usado en todo su
espectro y la fibra óptica. Y a diferencia de la DSN, \textbf{lo hizo a una
tasa prácticamente constante durante cuarenta años}: 0,179 décadas al año, es
decir un 51 por ciento anual, que es el valor que Jakob Nielsen postuló en
1998 como ley empírica y que su propia serie personal sigue cumpliendo dos
décadas después \parencite{nielsen2019}. Es, en mi opinión, el mejor ejemplo
disponible de la afirmación 3 en su forma completa, incluida la parte que la
DSN no cumple.
\end{quote}

\begin{table}[t]
\small
\caption{La velocidad de acceso doméstico, hito a hito. Las velocidades son
las nominales de la norma correspondiente en su sentido descendente. La
columna de la derecha es lo que cerró cada paradigma, y solo la primera es una
frontera demostrable. Normas: V.34 \parencite{ituv34}, V.90
\parencite{ituv90}, ADSL \parencite{itug9921}, ADSL2+ \parencite{itug9925} y
XGS-PON \parencite{itug98071}.}
\label{cuadro:internet}
\begin{tabularx}{\textwidth}{@{}p{2.15cm}>{\raggedright\arraybackslash}p{2.75cm}>{\raggedleft\arraybackslash}p{1.55cm}X@{}}
\toprule
\textbf{Paradigma} & \textbf{Norma y año} & \textbf{Velocidad} &
\textbf{Qué lo cerró} \\
\midrule
Canal de voz,\newline 3,1 kHz
& Bell 212A, 1984\newline ITU-T V.32bis, 1991\newline ITU-T V.34, 1994\newline
  ITU-T V.90, 1998
& 300 b/s\newline 14,4 kb/s\newline 33,6 kb/s\newline 56 kb/s
& \textbf{El teorema de Shannon.} Con 3,1 kHz de ancho y entre 30 y 35 dB de
  relación señal a ruido, la capacidad del canal está entre 31 y 36 kb/s:
  V.34 se sentó encima del límite. V.90 solo lo superó dejando de usar el
  canal analógico en un sentido \\
\addlinespace[2pt]
Par de cobre,\newline todo el espectro
& ITU-T G.992.1, 1999\newline ITU-T G.992.5, 2003\newline VDSL2 y DOCSIS 3.0,
  2006
& 1,5 Mb/s\newline 24 Mb/s\newline 100 Mb/s
& La atenuación del cobre, que crece con la frecuencia y con la longitud del
  bucle. Se compensó acercando la fibra al abonado, que ya es el paradigma
  siguiente disfrazado \\
\addlinespace[2pt]
Fibra hasta\newline el domicilio
& ITU-T G.984, 2003\newline ITU-T G.9807.1, 2016
& 1 Gb/s\newline 10 Gb/s
& \textbf{Nada físico.} La fibra tiene órdenes de magnitud de margen. Lo que
  frena hoy es la demanda: por encima de 1 Gb/s el usuario doméstico no
  distingue la diferencia \\
\bottomrule
\end{tabularx}
\end{table}

El Cuadro~\ref{cuadro:internet} y la Figura~\ref{fig:internet} desarrollan el
caso. Tres cosas merecen subrayarse.

\textbf{La primera asíntota es demostrable, igual que en la DSN.} La capacidad
de un canal de ancho $B$ con relación señal a ruido $S/N$ es

\begin{equation}
C \;=\; B \,\log_{2}\!\left(1 + \frac{S}{N}\right)
\label{eq:shannon}
\end{equation}

\parencite{shannon1948}. Para el canal telefónico, $B = 3100$ Hz y una
relación señal a ruido realista de entre 30 y 35 dB, la ecuación
\eqref{eq:shannon} da entre 30,9 y 36,0 kb/s. La norma V.34 alcanzó 33,6 kb/s
\parencite{ituv34}. No se quedó cerca del
límite: se sentó justo encima de él. Y la norma siguiente, V.90, no rompió esa
barrera mejorando el módem, porque eso es imposible: la rompió
\emph{cambiando el canal}, aprovechando que el extremo del proveedor ya estaba
digitalizado y que por tanto una de las dos direcciones se ahorraba una
conversión analógica \parencite{ituv90}. Eso es literalmente un cambio de
curva S forzado por una asíntota demostrada.

\textbf{La segunda: las tres eras tienen casi la misma pendiente.} El Panel B
de la Figura~\ref{fig:internet} las mide: 0,162 décadas al año en el canal de
voz, 0,177 en el cobre de banda ancha y 0,127 en la fibra, frente al 0,176 que
predice la regla del 50 por ciento anual de \textcite{nielsen2019}. Tres
medios físicos radicalmente distintos, tres industrias distintas, tres
regulaciones distintas, y la misma pendiente dentro del margen de una lectura
razonable. \textbf{Eso es la afirmación 3 cumplida entera}, y es lo que la DSN
no consigue.

\textbf{La tercera, y es la que menos esperaba: lo que está frenando la última
era no es la física.} La fibra tiene margen de sobra; XGS-PON entrega 10
gigabits y las normas de 25 y 50 gigabits ya están publicadas. Si la pendiente
del tramo de fibra baja a 0,127 décadas al año no es porque cueste más
conseguir la mejora, sino porque por encima de un gigabit el usuario doméstico
deja de notarla y deja de pagarla. \textbf{La restricción activa se ha movido
de la oferta a la demanda.} Es un desplazamiento distinto del de la DSN, donde
la restricción se movió de la física a las instituciones, pero el efecto sobre
la curva es el mismo, y ninguno de los dos aparece en el enunciado de las
cuatro afirmaciones.

\begin{figure}[!htb]
\centering
\begin{tikzpicture}

% ================= Panel A: la serie =====================================
\node[anchor=south west, font=\small\bfseries] at (-0.95,4.65)
     {Panel A. Velocidad de acceso doméstico a internet, 1984 a 2026};
\node[anchor=south west, font=\scriptsize\itshape, text=cGris] at (-0.95,4.35)
     {bits por segundo en sentido descendente, escala logarítmica};

\foreach \v/\r in {2/{100 b/s}, 3/{1 kb/s}, 4/{10 kb/s}, 5/{100 kb/s},
                   6/{1 Mb/s}, 7/{10 Mb/s}, 8/{100 Mb/s}, 9/{1 Gb/s},
                   10/{10 Gb/s}} {
  \draw[cGris!25] (0,{\ny{\v}}) -- (13.35,{\ny{\v}});
  \node[left, font=\scriptsize, text=cGris] at (-0.07,{\ny{\v}}) {\r};
}
\draw[cGris!70] (0,0) -- (0,{\ny{10.35}});
\draw[cGris!70] (0,0) -- (13.35,0);
\foreach \a in {1985,1990,1995,2000,2005,2010,2015,2020,2025} {
  \draw[cGris!70] ({\nx{\a}},0) -- ({\nx{\a}},-0.09);
  \node[below, font=\scriptsize, text=cGris] at ({\nx{\a}},-0.09) {\a};
}
\node[below, font=\scriptsize, text=cGris] at (6.6,-0.50) {año};

% recta de referencia de Nielsen: +50 % anual desde 300 b/s en 1984
\draw[cGris, dash pattern=on 2pt off 2pt, thick]
  ({\nx{1984}},{\ny{2.477}}) -- ({\nx{2026}},{\ny{9.875}});
\node[right, font=\scriptsize, text=cGris!20!black, rotate=31, inner sep=3pt]
     at ({\nx{2003}},{\ny{5.55}}) {ley de Nielsen: $+50$ \% anual};

% límite de Shannon del canal de voz
\fill[cLim!12] (0,{\ny{4.49}}) rectangle (13.35,{\ny{4.557}});
\draw[cLim, dashed] (0,{\ny{4.557}}) -- (13.35,{\ny{4.557}});
\node[left, font=\scriptsize, text=cLim, align=right, inner sep=3pt]
     at ({\nx{2026}},{\ny{3.72}})
     {límite de Shannon del canal\\telefónico de voz: 31 a 36 kb/s};

% era del canal de voz
\draw[cNet, very thick]
  ({\nx{1984}},{\ny{2.477}}) -- ({\nx{1988}},{\ny{2.477}}) --
  ({\nx{1988}},{\ny{3.380}}) -- ({\nx{1991}},{\ny{3.380}}) --
  ({\nx{1991}},{\ny{4.158}}) -- ({\nx{1994}},{\ny{4.158}}) --
  ({\nx{1994}},{\ny{4.459}}) -- ({\nx{1996}},{\ny{4.459}}) --
  ({\nx{1996}},{\ny{4.526}}) -- ({\nx{1998}},{\ny{4.526}}) --
  ({\nx{1998}},{\ny{4.748}}) -- ({\nx{2000}},{\ny{4.748}});
% era del cobre de banda ancha
\draw[cDato, very thick]
  ({\nx{2000}},{\ny{4.748}}) -- ({\nx{2000}},{\ny{6.176}}) --
  ({\nx{2003}},{\ny{6.176}}) -- ({\nx{2003}},{\ny{6.903}}) --
  ({\nx{2006}},{\ny{6.903}}) -- ({\nx{2006}},{\ny{7.380}}) --
  ({\nx{2010}},{\ny{7.380}}) -- ({\nx{2010}},{\ny{8.000}}) --
  ({\nx{2014}},{\ny{8.000}}) -- ({\nx{2014}},{\ny{8.477}}) --
  ({\nx{2016}},{\ny{8.477}});
% era de la fibra
\draw[cOpt, very thick]
  ({\nx{2016}},{\ny{8.477}}) -- ({\nx{2018}},{\ny{8.477}}) --
  ({\nx{2018}},{\ny{9.000}}) -- ({\nx{2022}},{\ny{9.000}}) --
  ({\nx{2022}},{\ny{9.301}}) -- ({\nx{2026}},{\ny{9.301}}) --
  ({\nx{2026}},{\ny{10.000}});

\node[left, font=\scriptsize, text=cNet, inner sep=3pt]
     at ({\nx{1999.6}},{\ny{5.35}}) {V.90};

% leyenda de los tres medios, al pie del panel
\draw[cNet, very thick] (0.15,-1.18) -- (0.85,-1.18);
\node[right, font=\scriptsize, text=cNet, inner sep=3pt] at (0.90,-1.18)
     {\textbf{canal de voz}: módem sobre 3,1 kHz};
\draw[cDato, very thick] (5.55,-1.18) -- (6.25,-1.18);
\node[right, font=\scriptsize, text=cDato, inner sep=3pt] at (6.30,-1.18)
     {\textbf{par de cobre}: ADSL, VDSL2, cable};
\draw[cOpt, very thick] (10.95,-1.18) -- (11.65,-1.18);
\node[right, font=\scriptsize, text=cOpt, inner sep=3pt] at (11.70,-1.18)
     {\textbf{fibra}: GPON, XGS-PON};

% ================= Panel B: la tasa por era ==============================
\begin{scope}[xshift=5.95cm, yshift=-6.30cm]
\node[anchor=south west, font=\small\bfseries] at (-4.75,3.42)
     {Panel B. La tasa de progreso, era por era, con la ley de Nielsen};

\barraU{2.35}{0.162}{cNet}{\textbf{Canal de voz}, 1984 a 1998\\2,27 décadas en
  14 años}{0,162}
\barraU{1.60}{0.177}{cDato}{\textbf{Cobre}, 2000 a 2012\\2,12 décadas en 12
  años}{0,177}
\barraU{0.85}{0.127}{cOpt}{\textbf{Fibra}, 2014 a 2026\\1,52 décadas en 12
  años}{0,127}
\barraU{0.10}{0.179}{cGris}{\textbf{Los 42 años juntos}\\7,52 décadas en 42
  años}{0,179}

\draw[cLim, thick] ({\ux{0.176}},-0.22) -- ({\ux{0.176}},2.80);
\node[above, font=\scriptsize, text=cLim, align=center, inner sep=2pt]
     at ({\ux{0.176}},2.80) {ley de Nielsen: 0,176};

\draw[cGris!70] (0,-0.26) -- ({\ux{0.195}},-0.26);
\foreach \t/\r in {0/0, 0.05/{0,05}, 0.10/{0,10}, 0.15/{0,15}} {
  \draw[cGris!70] ({\ux{\t}},-0.26) -- ({\ux{\t}},-0.35);
  \node[below, font=\scriptsize, text=cGris] at ({\ux{\t}},-0.35) {\r};
}
\node[below, font=\scriptsize, text=cGris] at ({\ux{0.09}},-0.72)
     {décadas de mejora por año};
\end{scope}

\end{tikzpicture}
\caption{Panel A: la misma figura de mérito y el mismo tipo de gráfico que la
Figura~\ref{fig:dsn}, aplicados al acceso doméstico a internet. Los tres
colores son tres medios físicos distintos. La recta gris es la regla del 50
por ciento anual de \textcite{nielsen2019}, trazada desde el primer punto y no
ajustada: la serie la sigue durante cuarenta años. La línea roja es la
capacidad de Shannon del canal telefónico, y se ve que la era del módem no se
detuvo por falta de ingenio sino porque chocó con ella. Panel B: las
pendientes de las tres eras son casi iguales, que es justamente lo que no
ocurre en la Red del Espacio Profundo. Elaboración propia sobre las normas
citadas en el Cuadro~\ref{cuadro:internet}.}
\label{fig:internet}
\end{figure}

\section{Las tres series juntas, y la respuesta que se sostiene}

La Figura~\ref{fig:sintesis} pone las tres series que he usado en la misma
unidad, décadas de mejora por año, y las tres se comportan de manera distinta:
la Red del Espacio Profundo decelera, el acceso a internet se mantiene y el
cómputo de entrenamiento de los modelos de aprendizaje automático acelera. Las
tres son tecnologías de información, las tres son funcionales en el sentido de
\textcite{koh2006} y las tres están medidas con figuras de mérito declaradas.

\begin{figure}[!htb]
\centering
\begin{tikzpicture}

% ================= Panel A: las tres series ==============================
\node[anchor=south west, font=\small\bfseries] at (-0.85,4.62)
     {Panel A. Tres tecnologías de información, tres comportamientos};
\node[anchor=south west, font=\scriptsize\itshape, text=cGris] at (-0.85,4.32)
     {décadas de mejora por año de la figura de mérito de cada una};

\foreach \v/\r in {0/0, 0.2/{0,2}, 0.4/{0,4}, 0.6/{0,6}, 0.8/{0,8}} {
  \draw[cGris!25] (0,{\ry{\v}}) -- (12.0,{\ry{\v}});
  \node[left, font=\scriptsize, text=cGris] at (-0.07,{\ry{\v}}) {\r};
}
\draw[cGris!70] (0,0) -- (0,{\ry{0.95}});
\draw[cGris!70] (0,0) -- (12.0,0);
\foreach \a in {1960,1970,1980,1990,2000,2010,2020,2030} {
  \draw[cGris!70] ({\rx{\a}},0) -- ({\rx{\a}},-0.09);
  \node[below, font=\scriptsize, text=cGris] at ({\rx{\a}},-0.09) {\a};
}
\node[below, font=\scriptsize, text=cGris] at (6.0,-0.50) {año};

% DSN
\draw[cX, very thick]
  ({\rx{1958}},{\ry{0.889}}) -- ({\rx{1976}},{\ry{0.889}}) --
  ({\rx{1976}},{\ry{0.076}}) -- ({\rx{1993}},{\ry{0.076}}) --
  ({\rx{1993}},{\ry{0.007}}) -- ({\rx{2005}},{\ry{0.007}}) --
  ({\rx{2005}},{\ry{0.121}}) -- ({\rx{2022}},{\ry{0.121}});
\draw[cX, very thick, dash pattern=on 4pt off 2pt]
  ({\rx{2022}},{\ry{0.121}}) -- ({\rx{2022}},{\ry{0.228}}) --
  ({\rx{2031}},{\ry{0.228}});
\node[right, font=\scriptsize\bfseries, text=cX, align=left, inner sep=3pt]
     at ({\rx{1978.5}},{\ry{0.80}}) {Red del Espacio Profundo:\\\emph{decelera}
     (afirmación 2)};
\draw[cX, -{Stealth[length=4pt]}, thin]
     ({\rx{1978}},{\ry{0.76}}) -- ({\rx{1973}},{\ry{0.58}});

% acceso a internet
\draw[cNet, very thick]
  ({\rx{1984}},{\ry{0.162}}) -- ({\rx{1998}},{\ry{0.162}}) --
  ({\rx{1998}},{\ry{0.177}}) -- ({\rx{2012}},{\ry{0.177}}) --
  ({\rx{2012}},{\ry{0.127}}) -- ({\rx{2026}},{\ry{0.127}});
\draw[cNet, -{Stealth[length=4pt]}, thin]
     ({\rx{1990}},{\ry{0.400}}) -- ({\rx{1990}},{\ry{0.205}});
\node[above, font=\scriptsize\bfseries, text=cNet, align=center,
      text width=3.45cm, inner sep=3pt]
     at ({\rx{1990}},{\ry{0.405}}) {acceso a internet:\\\emph{constante}
     (afirmación 3)};

% cómputo de aprendizaje automático
\draw[cKa, very thick]
  ({\rx{1990}},{\ry{0.170}}) -- ({\rx{2010}},{\ry{0.170}}) --
  ({\rx{2010}},{\ry{0.634}}) -- ({\rx{2022}},{\ry{0.634}});
\draw[cKa, -{Stealth[length=4pt]}, thin]
     ({\rx{2026}},{\ry{0.545}}) -- ({\rx{2022.5}},{\ry{0.620}});
\node[left, font=\scriptsize\bfseries, text=cKa, align=right,
      inner sep=3pt] at ({\rx{2030}},{\ry{0.475}})
     {cómputo de entrenamiento:\\\emph{acelera} (afirmación 4)};

% leyenda de las tres series, al pie del panel
\draw[cX, very thick] (0.15,-1.16) -- (0.85,-1.16);
\node[right, font=\scriptsize, text=cX, inner sep=3pt] at (0.90,-1.16)
     {\textbf{Red del Espacio Profundo}};
\draw[cNet, very thick] (5.05,-1.16) -- (5.75,-1.16);
\node[right, font=\scriptsize, text=cNet, inner sep=3pt] at (5.80,-1.16)
     {\textbf{acceso a internet}};
\draw[cKa, very thick] (8.70,-1.16) -- (9.40,-1.16);
\node[right, font=\scriptsize, text=cKa, inner sep=3pt] at (9.45,-1.16)
     {\textbf{cómputo de entrenamiento}};

\node[anchor=north west, font=\scriptsize\itshape, text=cGris!30!black,
      text width=11.6cm, align=left] at (-0.80,-1.56)
     {Las tres son tecnologías de información y las tres tienen una figura de
      mérito funcional bien declarada. Si las afirmaciones 2, 3 y 4 fueran
      leyes, estas tres líneas tendrían la misma forma. No la tienen, y por eso
      ninguna de las tres es una ley: son descripciones de casos.};

% ================= Panel B: los tres niveles =============================
\begin{scope}[yshift=-8.45cm]
\node[anchor=south west, font=\small\bfseries] at (-0.85,4.90)
     {Panel B. A qué altura es verdadera cada afirmación};

\draw[cOpt, thick, rounded corners=4pt] (0.10,0.10) rectangle (11.90,4.62);
\node[anchor=north west, font=\small\bfseries, text=cOpt] at (0.32,4.54)
     {FUNCIÓN: mover bits desde una sonda lejana};
\node[anchor=north west, font=\scriptsize, text=cOpt, text width=11.2cm]
     at (0.32,4.12)
     {\textbf{Afirmaciones 2 y 4:} son afirmaciones sobre la \emph{pendiente}
      de la envolvente funcional. No se deducen, se miden, y cada tecnología
      da un resultado distinto. Aquí: 20,8 décadas en 64 años, decelerando.};

\draw[cKa, thick, rounded corners=4pt] (0.45,0.22) rectangle (11.55,3.08);
\node[anchor=north west, font=\small\bfseries, text=cKa] at (0.67,3.00)
     {PARADIGMA: banda S, banda X, banda Ka, óptico};
\node[anchor=north west, font=\scriptsize, text=cKa, text width=10.5cm]
     at (0.67,2.58)
     {\textbf{Afirmación 3:} describe correctamente el mecanismo de relevo.
      Pero el hueco entre curvas, doce años entre 1993 y 2005, es parte del
      fenómeno y la afirmación no lo menciona.};

\draw[cLim, thick, rounded corners=4pt] (0.80,0.34) rectangle (11.20,1.52);
\node[anchor=north west, font=\scriptsize, text=cLim, text width=9.9cm]
     at (1.02,1.42)
     {\textbf{COMPONENTE} (máser, reflector, código): \textbf{afirmación 1},
      verdadera y aquí demostrable. Shannon, el fondo cósmico de microondas,
      la deformación del reflector, la ventana atmosférica.};
\end{scope}

\end{tikzpicture}
\caption{Panel A: las tres series en la misma unidad. La de la DSN es cálculo
propio sobre \textcite{descanso2015}; la de acceso a internet, cálculo propio
sobre las normas del Cuadro~\ref{cuadro:internet} y
\textcite{nielsen2019}; la del cómputo de entrenamiento procede de
\textcite{sevilla2022}, que mide duplicaciones cada 21,3 meses antes de 2010 y
cada 5,7 después. Panel B: la estructura de la respuesta. Las cuatro
afirmaciones no compiten: hablan de tres alturas distintas del mismo sistema.
Elaboración propia.}
\label{fig:sintesis}
\end{figure}

Eso es, para mí, la respuesta al planteamiento. \textbf{Las afirmaciones 1 y 3
son verdaderas porque describen mecanismos, y los mecanismos se pueden
verificar en cualquier caso.} La 1 describe qué le pasa a un componente y aquí
se puede demostrar con Shannon y con el fondo cósmico de microondas. La 3
describe cómo se relevan los paradigmas y aquí está dibujada en el propio
gráfico oficial de la JPL.

\textbf{Las afirmaciones 2 y 4 no son verdaderas ni falsas en general, porque
no son mecanismos sino mediciones}, y la medición da resultados opuestos según
la tecnología: la 2 se cumple en la DSN y no en el acceso a internet; la 4 no
se cumple en ninguna de las dos y sí en el cómputo de entrenamiento de
modelos. Presentarlas como leyes es el error; presentarlas como preguntas que
hay que contestar midiendo, para cada tecnología concreta y con una figura de
mérito declarada, es exactamente el método que propone \textcite{deweck2022} y
es lo único que sirve para decidir dónde invertir.

\section{Lo que me llevo del ejercicio}

\begin{enumerate}
\item \textbf{La altura a la que se mide decide la respuesta}, y por eso la
  primera obligación al discutir progreso tecnológico es declarar la figura de
  mérito y el nivel. Sin eso, las cuatro afirmaciones son indecidibles y la
  discusión se convierte en un intercambio de anécdotas.
\item \textbf{La media de sesenta años no describe ningún año.} La DSN
  promedia 0,325 décadas anuales y no ha tenido jamás un año normal a esa tasa:
  la cifra la produce un tramo inicial de dieciocho años que nunca se repitió.
  Extrapolar la media histórica de una tecnología madura es el error de
  previsión más caro y el más frecuente.
\item \textbf{El hueco entre curvas S importa tanto como el salto.} Doce años
  planos entre la banda X y la banda Ka. Para un director de tecnología, la
  variable de gestión no es cuándo llega la curva siguiente sino cuánto dura el
  hueco, y eso depende de decisiones de capital tomadas con años de antelación.
\item \textbf{Conviene saber si la asíntota es demostrable o solo aparente.}
  Shannon y el fondo cósmico son topes que no se negocian. La lentitud de la
  fibra hoy, en cambio, es demanda insuficiente, y eso se revierte en cuanto
  aparece una aplicación que la justifique. Tratar las dos igual lleva a
  abandonar tecnologías que aún tenían recorrido.
\item \textbf{La restricción activa se desplaza, y casi nunca avisa.} En la DSN
  pasó de la física a la institución; en el acceso a internet, de la oferta a
  la demanda. Ninguno de los dos desplazamientos aparece en el enunciado de las
  cuatro afirmaciones, y los dos son la información que de verdad hace falta
  para decidir.
\end{enumerate}

\section{Limitaciones}

\begin{enumerate}
\item \textbf{Los datos de la DSN son lecturas visuales de un gráfico
  logarítmico}, no una tabla publicada. Con una tolerancia de $\pm$0,2 décadas
  por punto, las pendientes de las eras de banda X y Ka, que se calculan sobre
  recorridos de 1,2 y 1,9 décadas, arrastran un error relativo que puede llegar
  al 20 por ciento. El orden de magnitud de la conclusión, que la primera era
  fue siete veces más rápida que la última, resiste ese error; el valor exacto
  de cada pendiente, no.
\item \textbf{Las fronteras entre eras las he elegido yo.} Situar el fin de la
  era de banda X en 1993 y el comienzo de la de Ka en 2006 es una decisión
  defendible pero no la única. Moverlas dos o tres años cambia las pendientes,
  aunque no el sentido de la comparación.
\item \textbf{Cito la entrevista al doctor Deutsch de memoria}, sin
  transcripción ni marca de tiempo, y por tanto la atribución de la idea de que
  las limitaciones actuales son más institucionales que físicas es mi lectura y
  puede no recoger fielmente lo que dijo. Los argumentos que construyo sobre
  ella, la meseta de doce años y el desplazamiento de la restricción activa,
  se sostienen sobre los datos del gráfico y no sobre la entrevista.
\item \textbf{La serie de velocidad de acceso a internet es una reconstrucción
  a partir de las normas}, no una medición de velocidades reales entregadas al
  abonado, que siempre son menores. Usar la velocidad nominal de la norma es
  coherente a lo largo de los 42 años, que es lo que exige el cálculo de una
  pendiente, pero sobrestima el nivel absoluto en todos los puntos.
\item \textbf{El reparto entre ganancia extraída y margen restante del
  Cuadro~\ref{cuadro:asintotas} mezcla unidades cómodamente.} Convertir
  decibelios de margen de codificación en décadas de tasa de datos supone que
  la relación es lineal, y solo lo es aproximadamente y en un régimen
  determinado. El porcentaje de consumo es una buena guía y no una medida
  exacta.
\item \textbf{La tercera era del acceso a internet aún no ha terminado}, de
  modo que su pendiente de 0,127 décadas al año es provisional: si las normas
  de 25 y 50 gigabits se despliegan a escala en los próximos años, ese número
  subirá y mi conclusión sobre el freno por demanda quedaría debilitada.
\item \textbf{No he comprobado personalmente ninguna de las tres series.} Todas
  se apoyan en fuentes secundarias, aunque las tres sean fuentes primarias de
  sus respectivos dominios. En particular, los tiempos de duplicación del
  cómputo de aprendizaje automático proceden de un único trabajo
  \parencite{sevilla2022} y su régimen más reciente es corto y sensible a la
  elección de la muestra.
\end{enumerate}

\printbibliography

\end{document}
